\documentclass[12 pt]{article}
\usepackage[utf8]{inputenc}
\usepackage{matlab-prettifier}
\usepackage[portuguese]{babel}
\usepackage{indentfirst}
\usepackage{graphicx}
\usepackage{float}
\usepackage{subcaption}
\usepackage[font=small,labelfont=bf]{caption}
\definecolor{mygreen}{RGB}{28,172,0} % color values Red, Green, Blue
\definecolor{myyellow}{rgb}{1.0, 1.0, 0.8}
\usepackage{mathtools}
\usepackage{multirow}
\usepackage{comment}
\usepackage{xcolor}
\usepackage{colortbl}
\usepackage[normalem]{ulem}               % to striketrhourhg text
\usepackage{amsmath}
\usepackage{amsfonts}
\usepackage{hyperref}
\usepackage{tcolorbox}
\tcbset{before upper={\setlength{\parindent}{2em}}}
\usepackage{longtable}
\usepackage{enumitem}

\newcommand\redout{\bgroup\markoverwith
{\textcolor{red}{\rule[0.5ex]{2pt}{0.8pt}}}\ULon}
\renewcommand{\lstlistingname}{Código}% Listing -> Algorithm
\renewcommand{\lstlistlistingname}{Lista de \lstlistingname s}% List of Listings -> List of Algorithms

\usepackage[top=3cm,left=2cm,bottom=2cm, right=2cm]{geometry}
\usepackage{tikz}
\usetikzlibrary{decorations.pathreplacing}
\usetikzlibrary{automata}
\usetikzlibrary{positioning}
\usetikzlibrary{arrows.meta, positioning}

\usepackage{adjustbox}

\usepackage{mdframed}
\newmdenv[
    backgroundcolor=gray!15,       % fundo cinza claro
    linecolor=gray!60,             % borda cinza mais escura
    linewidth=1pt,
    roundcorner=5pt,
    skipabove=10pt,
    skipbelow=10pt,
    innertopmargin=1em,
    innerbottommargin=1em,
    innerleftmargin=1em,
    innerrightmargin=1em,
    splittopskip=\topskip,
    firstextra=\setlength{\parindent}{2em}, 
]{respostabox@inner}

\newenvironment{resposta}
    {\begin{respostabox@inner}
    \setlength{\parindent}{2em}~
    }
  {\end{respostabox@inner}}

% Configuração para destacar a sintaxe do Python
\lstset{ 
    language=Python,                     % A linguagem do código
    backgroundcolor=\color{myyellow}, % A cor do fundo 
    basicstyle=\ttfamily\footnotesize,   % O estilo do texto básico
    keywordstyle=\color{blue},           % Cor das palavras-chave
    stringstyle=\color{red},             % Cor das strings
    commentstyle=\color{mygreen},          % Cor dos comentários
    numbers=left,                        % Números das linhas à esquerda
    numberstyle=\tiny\color{gray},       % Estilo dos números das linhas
    stepnumber=1,                        % Número de linhas entre os números das linhas
    frame=single,                        % Moldura ao redor do código
    breaklines=true,                     % Quebra automática das linhas longas
    captionpos=t,                        % Posição da legenda
    showstringspaces=false               % Não mostra espaços em branco nas strings
    extendedchars=true,
    literate={º}{{${ }^{\underline{o}}$}}1 {á}{{\'a}}1 {à}{{\`a}}1 {ã}{{\~a}}1 {é}{{\'e}}1 {É}{{\'E}}1 {ê}{{\^e}}1 {ë}{{\"e}}1 {í}{{\'i}}1 {ç}{{\c{c}}}1 {Ç}{{\c{C}}}1 {õ}{{\~o}}1 {ó}{{\'o}}1 {ô}{{\^o}}1 {ú}{{\'u}}1 {â}{{\^a}}1 {~}{{$\sim$}}1
}


\title{%
\textbf{\huge Universidade Federal do Rio de Janeiro} \par
\textbf{\LARGE Instituto Alberto Luiz Coimbra de Pós-Graduação e Pesquisa de Engenharia} \par

\includegraphics[width=8cm]{COPPE UFRJ.png} \par

\textbf{Programa de Engenharia de Sistemas e Computação} \par

CPS767 - Algoritmos de Monte Carlo e Cadeias de Markov  \par

Prof. Daniel Ratton Figueiredo\par

\vspace{1\baselineskip}
\textbf{\textit{3ª Lista de Exercícios}} \par
}

\author{Luiz Henrique Souza Caldas\\email: lhscaldas@cos.ufrj.br}

\date{\today}

\begin{document}
\maketitle

% \tableofcontents

\section*{Questão 1: Passeios aleatórios enviesados}
Considere um grafo não direcionado $G = (V,E)$ com peso nas arestas, tal que $w_{ij} > 0$ para toda aresta $(i, j) \in E$. Considere um andarilho aleatório que caminha por este grafo em tempo discreto, mas cujos passos são enviesados pelos pesos das arestas. Em particular, a probabilidade do andarilho ir do vértice $i$ para o vértice $j$ é dada por $w_{ij}/W_i$, onde $W_i = \sum_j w_{ij}$ (soma dos pesos das arestas incidentes ao vértice $i \in V$). Temos assim um passeio aleatório enviesado linearmente pelos pesos das arestas.

\begin{enumerate}
    \item Mostre que este passeio aleatório induz uma cadeia de Markov calculando a matriz de transição de probabilidade.
    \begin{resposta}
        
    Conforme o enunciado, temos que a probabilidade de transição do estado $i$ para o estado $j$ é dada por:
    $$ P_{ij} = \frac{w_{ij}}{W_i}, $$
    onde $W_i = \sum_{j} w_{ij}$ é a soma dos pesos das arestas incidentes ao vértice $i$ e $w_{ij}>0$ é o peso da aresta que liga os vértices $i$ e $j$.

    Como $w_{ij} > 0$ para todas as arestas $(i,j) \in E$ e $W_i > 0$, temos que $P_{ij} \geq 0$, sendo estritamente positivo sempre que houver uma aresta entre $i$ e $j$.

    
    Além disso,

    $$ \sum_{j} P_{ij} = \sum_{j} \frac{w_{ij}}{W_i} = \frac{1}{W_i} \sum_{j} w_{ij} = \frac{W_i}{W_i} = 1 $$

    Assim, temos que a soma das probabilidades de transição do estado $i$ para todos os outros estados $j$ é igual a 1, o que caracteriza uma matriz de transição de probabilidade. 
    
    Portanto, \underline{o passeio aleatório enviesado induz uma cadeia de Markov}.
        
    \end{resposta}
    \item Determine a distribuição estacionária desta cadeia de Markov (dica: use o método da inspeção).
    \begin{resposta}
        Uma distribuição estacionária $\pi$ é uma distribuição de probabilidade que satisfaz a seguinte equação:
        $$ \pi P = \pi $$

        Para cada estado $i$ da distribuição estacionária, temos que:
        $$ \pi_i = \sum_{j} \pi_j P_{ji} $$

        Usando o resultado do item anterior, temos que:
        $$ \pi_i = \sum_{j} \pi_j \frac{w_{ji}}{W_j} $$

        Supondo que a distribuição estacionária é proporcional à soma do peso das arestas incidentes ao vértice $j$, ou seja, $\pi_j = Z W_j$, onde $Z>0$ é uma constante, temos que:
        $$ \pi_i = \sum_{j} (Z W_j) \frac{w_{ji}}{W_j} = Z \sum_{j} w_{ji}$$

        Como o grafo é não direcionado, temos que $w_{ij} = w_{ji}$, ou seja:
        $$ \pi_i = Z \sum_{j} w_{ji} = Z \sum_{j} w_{ij} = Z W_i $$

        Porém, sabemos que $\sum_{k} \pi_k = 1$, ou seja:
        $$ \sum_{k} \pi_k = \sum_{k} Z W_k = 1 \quad \Rightarrow \quad Z = \frac{1}{\sum_{k} W_k}$$
    
        Portanto, a distribuição estacionária é dada por:
        $$ \boxed{\pi_i = \frac{W_i}{\sum_{k} W_k}} $$


    \end{resposta}
    \item Determine se esta cadeia de Markov é reversível no tempo.
    \begin{resposta}
        Uma Cadeia de Markov é dita reversível no tempo se, e somente se, satizfaz a equação de fluxo de massa de probabilidade:
        $$ \pi_i P_{ij} = \pi_j P_{ji} $$

        Aplicando a distribuição estacionária que encontramos no item anterior, temos que:
        $$ \pi_i P_{ij} = \frac{W_i}{\sum_{k} W_k} \cdot \frac{w_{ij}}{W_i} = \frac{w_{ij}}{\sum_{k} W_k}  \quad \text{e} \quad \pi_j P_{ji} = \frac{W_j}{\sum_{k} W_k} \cdot \frac{w_{ji}}{W_j} = \frac{w_{ji}}{\sum_{k} W_k}$$
        Como $w_{ij} = w_{ji}$, temos que:
        $$ \boxed{\pi_i P_{ij} = \pi_j P_{ji}} $$
        Portanto, \underline{a cadeia de Markov é reversível} no tempo.
    \end{resposta}
\end{enumerate}

\section*{Questão 2: Amostras de Modelos de Mistura}

Considere a seguinte função de probabilidade:
\[
p(x) = \alpha p_B(x; n, p_1) + (1 - \alpha)p_B(x; n, p_2),
\]
onde $p_B(x; n, p)$ é a probabilidade associada ao valor $x$ da binomial com parâmetros $n$ e $p$, e $\alpha \in [0,1]$ é um peso. Trata-se de um modelo de mistura de duas binomiais com diferentes valores de $p$, com pesos dados por $\alpha$ e $1 - \alpha$. Considere duas variáveis aleatórias $X$ e $K$, representando o valor de $X \in [0, n]$ e a binomial utilizada $K \in \{1, 2\}$. Queremos gerar amostras de acordo com $p(x)$.

\begin{itemize}
  \item Determine as distribuições de probabilidade condicionais $P(X|K)$ e $P(K|X)$. Dica: utilize a regra de Bayes no segundo caso.
  \begin{resposta}
    A distribuição de $X$ dado $K$ é simplesmente uma binomial com parâmetro $p_k$, para $k \in \{1, 2\}$:
    $$ P (X=x|K=k) = p_B(x; n, p_k)$$
    $$\boxed{P (X=x|K=k) = \binom{n}{x} p_k^x (1 - p_k)^{n - x}} $$
    
    A distribuição de $K$ dado $X$ é obtida pela regra de Bayes:
    $$ P (K=k|X=x) = \frac{P (X=x|K=k) \cdot P(K=k)}{p(x)} $$

    Como $P(K = 1) = \alpha$ e $P(K = 2) = 1 - \alpha$, temos:
    $$ \boxed{P (K=k|X=x) = \frac{p_B(x; n, p_k) \cdot p(K=k)}{\alpha p_B(x; n, p_1) + (1 - \alpha)p_B(x; n, p_2)}} $$
  \end{resposta}
  \item Determine a distribuição de probabilidade conjunta $P(X, K)$.
  \begin{resposta}
    A distribuição conjunta é dada pela regra do produto:
    $$ P(X=x, K=k) = P(X=x|K=k) \cdot P(K=k) $$

    Como $ P (X=x|K=k) = p_B(x; n, p_k)$,  $P(K = 1) = \alpha$, $P(K = 2) = 1 - \alpha$, temos:
    $$ \boxed{P(X=x, K=k) =p_B(x; n, p_k)  \cdot P(K=k)} $$
  \end{resposta}
  \item Utilize a técnica de Gibbs Sampling para gerar amostras de $X$. Mostre como construir a cadeia de Markov e determine a transição entre os estados.
\begin{resposta}
  Cada estado da cadeia de Markov é representado por um par $V = (X, K)$, onde $X \in [0, n]$ é o valor da variável aleatória e $K \in \{1, 2\}$ indica qual das duas binomiais foi utilizada.

  As transições entre os estados ocorrem por meio da atualização de uma única variável por vez, mantendo a outra fixa:

  \begin{itemize}
    \item Para transições em que apenas $X$ é atualizado (com $K$ fixo), a probabilidade de transição para o novo estado $(X', K)$ é dada por:
    $$
    P_{(X, K), (X', K)} = \frac{1}{2} P[X'|K] = \frac{1}{2} \cdot p_B(X'; n, p_K)
    $$

    \item Para transições em que apenas $K$ é atualizado (com $X$ fixo), a probabilidade de transição para o novo estado $(X, K')$ é dada por:
    $$
    P_{(X, K), (X, K')} = \frac{1}{2} P[K'|X] = \frac{1}{2} \cdot \frac{p_B(X; n, p_{K'}) \cdot P(K')}{\alpha \cdot p_B(X; n, p_1) + (1 - \alpha) \cdot p_B(X; n, p_2)}
    $$
    onde $P(K' = 1) = \alpha$ e $P(K' = 2) = 1 - \alpha$.
  \end{itemize}

  Dessa forma, a cadeia se move no espaço de pares $(X, K)$, com transições definidas pelas distribuições condicionais do modelo de mistura. A estrutura garante que a cadeia seja reversível e tenha como distribuição estacionária a distribuição conjunta $P(X, K)$.
\end{resposta}

  \newpage
  \item Para $n = 2$, $p_1 = 0{,}1$, $p_2 = 0{,}8$, $\alpha = 0{,}3$, desenhe a cadeia de Markov com todas as transições.
  \begin{resposta}

  \begin{center}
\renewcommand{\arraystretch}{1.3}
\begin{tabular}{|l|c|c|}
  \hline
  Probabilidade & Valor & Valor/2 \\
  \hline
  $P(X=0 \mid K=1)$ & $0.8100$ & $0.4050$ \\
  $P(X=1 \mid K=1)$ & $0.1800$ & $0.0900$ \\
  $P(X=2 \mid K=1)$ & $0.0100$ & $0.0050$ \\
  $P(X=0 \mid K=2)$ & $0.0400$ & $0.0200$ \\
  $P(X=1 \mid K=2)$ & $0.3200$ & $0.1600$ \\
  $P(X=2 \mid K=2)$ & $0.6400$ & $0.3200$ \\
  \hline
  $P(K=1 \mid X=0)$ & $0.8967$ & $0.4484$ \\
  $P(K=2 \mid X=0)$ & $0.1033$ & $0.0517$ \\
  $P(K=1 \mid X=1)$ & $0.1942$ & $0.0971$ \\
  $P(K=2 \mid X=1)$ & $0.8058$ & $0.4029$ \\
  $P(K=1 \mid X=2)$ & $0.0067$ & $0.0034$ \\
  $P(K=2 \mid X=2)$ & $0.9933$ & $0.4967$ \\
  \hline
\end{tabular}
\end{center}


\begin{center}
\begin{tikzpicture}[->, >=Stealth, node distance=3cm,
    every state/.style={circle, draw, minimum size=1.6cm},
    shorten >=1pt]

  % Estados (dispostos em grade 3x2)
  \node[state] (X0K1) at (0, 0) {$(0,1)$};
  \node[state] (X1K1) at (3, 0) {$(1,1)$};
  \node[state] (X2K1) at (6, 0) {$(2,1)$};

  \node[state] (X0K2) at (0, -4) {$(0,2)$};
  \node[state] (X1K2) at (3, -4) {$(1,2)$};
  \node[state] (X2K2) at (6, -4) {$(2,2)$};

  % Transições X | K=1 (p=0.1)
  \draw[bend left=10] (X0K1) to node[above] {\scriptsize $0.0900$} (X1K1);
  \draw[bend left=30] (X0K1) to node[above] {\scriptsize $0.0050$} (X2K1);
  \draw[bend left=10] (X1K1) to node[below] {\scriptsize $0.4050$} (X0K1);
  \draw[bend left=10] (X1K1) to node[above] {\scriptsize $0.0050$} (X2K1);
  \draw[bend left=10] (X2K1) to node[below] {\scriptsize $0.0900$} (X1K1);
  \draw[bend left=30] (X2K1) to node[below] {\scriptsize $0.4050$} (X0K1);

  % Transições X | K=2 (p=0.8)
  \draw[bend left=10] (X0K2) to node[above] {\scriptsize $0.1600$} (X1K2);
  \draw[bend left=30] (X0K2) to node[above] {\scriptsize $0.3200$} (X2K2);
  \draw[bend left=10] (X1K2) to node[below] {\scriptsize $0.0200$} (X0K2);
  \draw[bend left=10] (X1K2) to node[above] {\scriptsize $0.3200$} (X2K2);
  \draw[bend left=10] (X2K2) to node[below] {\scriptsize $0.1600$} (X1K2);
  \draw[bend left=30] (X2K2) to node[below] {\scriptsize $0.0200$} (X0K2);

  % Transições K | X fixo
  \draw[bend left=10] (X0K1) to node[right] {\scriptsize $0.0517$} (X0K2);
  \draw[bend left=10] (X0K2) to node[left] {\scriptsize $0.4484$} (X0K1);

  \draw[bend left=20] (X1K1) to node[right] {\scriptsize $0.4029$} (X1K2);
  \draw[bend left=20] (X1K2) to node[left] {\scriptsize $0.0971$} (X1K1);

  \draw[bend left=10] (X2K1) to node[right] {\scriptsize $0.4967$} (X2K2);
  \draw[bend left=10] (X2K2) to node[left] {\scriptsize $0.0034$} (X2K1);

  % Self-loops
  \draw[loop above] (X0K1) to node[above] {\scriptsize $0.8534$} (X0K1);
  \draw[->] (X1K1) to [out=120, in=60, looseness=5] node[above] {\scriptsize $0.1871$} (X1K1);
  \draw[loop above] (X2K1) to node[above] {\scriptsize $0.0084$} (X2K1);
  \draw[loop below] (X0K2) to node[below] {\scriptsize $0.0717$} (X0K2);
  \draw[->] (X1K2) to [out=-120, in=-60, looseness=5] node[below] {\scriptsize $0.5629$} (X1K2);
  \draw[loop below] (X2K2) to node[below] {\scriptsize $0.8167$} (X2K2);


\end{tikzpicture}
\end{center}
OBS: os valores das probabilidades dos selfloops são a soma das probabilidades de transição mantendo $X$ fixo e mantendo $K$ fixo, ou seja, $P(X=x|K=k) + P(K=k|X=x)$.

  \end{resposta}
  \item Descreva como utilizar a cadeia de Markov para gerar amostras.
  \begin{resposta}

    \begin{enumerate}
  \item Inicializar a cadeia em um estado arbitrário $(X, K)$, por exemplo $(0, 1)$.
  
  \item Escolher aleatoriamente qual variável será atualizada:
  \begin{itemize}
    \item Com probabilidade $\frac{1}{2}$, atualizar $X$;
    \item Com probabilidade $\frac{1}{2}$, atualizar $K$.
  \end{itemize}

  \item Atualizar a variável escolhida:
  \begin{itemize}
    \item Se for $X$, sorteie $X'$ da distribuição binomial $P(X \mid K)$ com parâmetro $p_K$;
    \item Se for $K$, sorteie $K'$ da distribuição $P(K \mid X)$, calculada pela regra de Bayes.
  \end{itemize}
  A variável não escolhida permanece inalterada.

  \item Atualizar o estado da cadeia para o novo par $(X', K')$.

  \item Repitir o processo. Após $\tau_{\epsilon}$ passos (tempo de mistura), o valor de $X$ pode ser considerado uma amostra da distribuição desejada.

  \item Repitir os passos acima para obter quantas amostras forem necessárias, aguardando sempre pelo menos $\tau_{\epsilon}$ passos entre duas amostras consecutivas.
\end{enumerate}

O tempo de mistura $\tau_{\epsilon}$ pode ser limitado superiormente por:
$$ \tau_{\epsilon} \leq \frac{\log{1/(\pi_o \epsilon)}}{\delta}$$
onde $\pi_o$ é a probabilidade do estado menos provável na distribuição estacionária da cadeia de Markov, $\epsilon$ é a tolerância desejada e $\delta$ é o vão espectral, definido por:
$$ \delta = 1 - |\lambda_2|$$
onde $\lambda_2$ é o segundo maior autovalor da matriz de transição da cadeia de Markov.
  \end{resposta}
\end{itemize}


\section*{Questão 3: Tempo de mistura}
Considere um processo estocástico que inicia no estado $1$ e a cada instante de tempo incrementa o valor do estado em uma unidade com probabilidade $p$ ou retorna ao estado inicial (estado $1$) com probabilidade $1-p$. No estado $n$ o processo não cresce mais, e se mantém neste estado com probabilidade $p$. Assuma que $n = 10$ e que $p \in \{0.25, 0.5, 0.75\}$.

\begin{enumerate}
    \item Construa a cadeia de Markov deste processo mostrando a matriz de transição de probabilidade em função de $p$.
    \begin{resposta}

        \begin{tikzpicture}[->, >=stealth, node distance=2.2cm,
            every state/.style={circle, draw, minimum size=1.2cm}]
            % States
            \node[state] (1) {1};
            \node[state] (2) [right=of 1] {2};
            \node[state] (3) [right=of 2] {3};
            \node (dots) [right=of 3] {\(\cdots\)};
            \node[state] (10) [right=of dots] {10};

            % Forward transitions
            \draw (1) -- node[above] {\(p\)} (2);
            \draw (2) -- node[above] {\(p\)} (3);
            \draw (3) -- node[above] {\(p\)} (dots);
            \draw (dots) -- node[above] {\(p\)} (10);

            % Return transitions
            \draw[bend left=30] (2) to node[below] {\(1 - p\)} (1);
            \draw[bend left=35] (3) to node[below] {\(1 - p\)} (1);
            \draw[bend left=40] (dots) to node[below] {\(1 - p\)} (1);
            \draw[bend left=45] (10) to node[below] {\(1 - p\)} (1);

            % Self-loops
            \draw (1) edge[loop above] node {\(1 - p\)} (1);
            \draw (10) edge[loop above] node {\(p\)} (10);

        \end{tikzpicture}

        $$
        P_{ij} = \begin{cases}
            1-p & \text{se } j = 1 \\
            p & \text{se } (j = i + 1) \lor (i=j=10)\\
            0 & \text{caso contrário}
        \end{cases}
        $$

        Matriz de transição (parcial):
                $$
                \begin{bmatrix}
                p & 1-p & 0    & \cdots & 0 \\
                p & 0    & 1-p &        & 0 \\
                \vdots &     &     & \ddots & \vdots \\
                p & 0    & 0    & \cdots & 1-p
                \end{bmatrix}
                $$

    \end{resposta}
    \item Determine numericamente o vão espectral da cadeia de Markov para cada valor de $p$.
    \begin{resposta}

        O vão espectral é dado por:
        $$\delta = 1 - |\lambda_2| $$
        
            onde $\lambda_2$ é o segundo maior autovalor (em módulo) da matriz de transição $P$.

        Utilizando um código em Python, com a biblioteca NumPy, podemos calcular os autovalores da matriz de transição $P$ para cada valor de $p$. O código pode ser visto no repositório do Github cujo link encontra-se no fim deste relatório.

        \begin{itemize}        
            \item \textbf{Para $p = 0{,}25$:}
            \begin{itemize}
                \item Matriz de transição (parcial):
                $$
                \begin{bmatrix}
                0.75 & 0.25 & 0    & \cdots & 0 \\
                0.75 & 0    & 0.25 &        & 0 \\
                \vdots &     &     & \ddots & \vdots \\
                0.75 & 0    & 0    & \cdots & 0.25
                \end{bmatrix}
                $$
                \item Maiores autovalores em módulo:
                $$ |\lambda_1| = 1.0, \quad |\lambda_2| \approx 1.7125 \times 10^{-4} $$
                \item Vão espectral:
                $$ \boxed{\delta = 0.9998287} $$
            \end{itemize}
        
            \item \textbf{Para $p = 0{,}5$:}
            \begin{itemize}
                \item Matriz de transição (parcial):
                $$
                \begin{bmatrix}
                0.5 & 0.5 & 0    & \cdots & 0 \\
                0.5 & 0   & 0.5  &        & 0 \\
                \vdots &     &     & \ddots & \vdots \\
                0.5 & 0   & 0    & \cdots & 0.5
                \end{bmatrix}
                $$
                \item Maiores autovalores em módulo:
                $$ |\lambda_1| = 1.0, \quad |\lambda_2| \approx 3.4251 \times 10^{-4} $$
                \item Vão espectral:
                $$ \boxed{\delta = 0.9996575} $$
            \end{itemize}
        
            \item \textbf{Para $p = 0{,}75$:}
            \begin{itemize}
                \item Matriz de transição (parcial):
                $$
                \begin{bmatrix}
                0.25 & 0.75 & 0    & \cdots & 0 \\
                0.25 & 0    & 0.75 &        & 0 \\
                \vdots &     &     & \ddots & \vdots \\
                0.25 & 0    & 0    & \cdots & 0.75
                \end{bmatrix}
                $$
                \item Maiores autovalores em módulo:
                $$ |\lambda_1| = 1.0, \quad |\lambda_2| \approx 4.3546 \times 10^{-6} $$
                \item Vão espectral:
                $$ \boxed{\delta = 0.9999956} $$
            \end{itemize}
        \end{itemize}
        


    \end{resposta}
    \item Determine numericamente a distribuição estacionária para cada valor de $p$, e indique o estado de menor probabilidade.
    \begin{resposta}
        O estado estacionário $\pi$ de uma cadeia de Markov é um vetor de probabilidade tal que:

        $$ \pi P = \pi $$

        Ou seja, ele é um autovetor à esquerda da matriz de transição $P$ associado ao autovalor $\lambda = 1$. 
        A distribuição estacionária pode ser obtida utilizando novamente um código Python com a biblioteca NumPy para computar os autovalores e autovetores de $P^\top$, selecionando aquele correspondente a $\lambda = 1$ e normalizando-o para que $\sum_i \pi_i = 1$.

        \begin{itemize}
            \item \textbf{Para $p = 0{,}25$:}
            \begin{itemize}
                \item Distribuição estacionária:
                $$
                \pi = 
                \begin{bmatrix}
                0.75 \\
                0.1875 \\
                0.046875 \\
                0.01171875 \\
                0.00292969 \\
                0.00073242 \\
                0.00018311 \\
                0.00004578 \\
                0.00001144 \\
                \boxed{0.00000381}
                \end{bmatrix}
                $$
                \item Estado de menor probabilidade: 10, com probabilidade $\pi_o = 0.00000381$.
            \end{itemize}

            \newpage
            \item \textbf{Para $p = 0{,}5$:}
            \begin{itemize}
                \item Distribuição estacionária:
                $$
                \pi = 
                \begin{bmatrix}
                0.5 \\
                0.25 \\
                0.125 \\
                0.0625 \\
                0.03125 \\
                0.015625 \\
                0.0078125 \\
                0.00390625 \\
                \boxed{0.00195312} \\
                \boxed{0.00195312}
                \end{bmatrix}
                $$
                \item Estado de menor probabilidade: 9 e 10, com probabilidade $\pi_o = 0.00195312$.
            \end{itemize}

            \newpage
            \item \textbf{Para $p = 0{,}75$:}
            \begin{itemize}
                \item Distribuição estacionária:
                $$
                \pi = 
                \begin{bmatrix}
                0.25 \\
                0.1875 \\
                0.140625 \\
                0.10546875 \\
                0.07910156 \\
                0.05932617 \\
                0.04449463 \\
                0.03337097 \\
                \boxed{0.02502823} \\
                0.07508469
                \end{bmatrix}
                $$
                \item Estado de menor probabilidade: 9, com probabilidade $\pi_o = 0.02502823$.
            \end{itemize}
        \end{itemize}



    \end{resposta}
    \item Utilizando os dados acima, determine um limitante inferior e superior para o tempo de mistura quando $\epsilon = 10^{-6}$ para cada valor de $p$.
    \begin{resposta}
        O tempo de mistura $\tau_\epsilon$ representa, de forma aproximada, o número de passos necessários para que a distribuição da cadeia de Markov esteja a uma distância total menor que $\epsilon$ da distribuição estacionária, independentemente do estado inicial.
        
        Ele pode ser estimado por:
        
        $$ \left( \frac{1}{\delta} - 1 \right) \log \frac{1}{2\epsilon} \leq \tau_\epsilon \leq \frac{1}{\delta} \log \frac{1}{\pi_o \epsilon}, $$
        
        onde $\delta$ é o vão espectral da matriz de transição, e $\pi_o = \min_i \pi_i$ é a menor entrada da distribuição estacionária.
        
        \begin{itemize}
            \item \textbf{Para $p = 0{,}25$:}
            \begin{itemize}
                \item Limite inferior estimado:
                $$ \tau_\epsilon \geq 0.0022476286040300246 $$
                \item Limite superior estimado:
                $$ \tau_\epsilon \leq 26.296663189659515 $$
            \end{itemize}
        
            \item \textbf{Para $p = 0{,}5$:}
            \begin{itemize}
                \item Limite inferior estimado:
                $$ \tau_\epsilon \geq 0.004496027298161922 $$
                \item Limite superior estimado:
                $$ \tau_\epsilon \leq 20.060706093968722 $$
            \end{itemize}
        
            \item \textbf{Para $p = 0{,}75$:}
            \begin{itemize}
                \item Limite inferior estimado:
                $$ \tau_\epsilon \geq 5.7142421447070306 \times 10^{-5} $$
                \item Limite superior estimado:
                $$ \tau_\epsilon \leq 17.50333771810467 $$
            \end{itemize}
        \end{itemize}
    \end{resposta}

    \item O que você pode concluir sobre a influência de $p$ no tempo de mistura?
    \begin{resposta}
        À medida que o valor de $p$ aumenta, o tempo de mistura da cadeia de Markov tende a diminuir. Isso ocorre porque valores maiores de $p$ favorecem a progressão pelos estados, permitindo que o processo alcance diferentes regiões do espaço de estados mais rapidamente antes de retornar ao início. Como consequência, a distribuição de probabilidade se aproxima da distribuição estacionária com maior eficiência. Por outro lado, quando $p$ é pequeno, a cadeia retorna com mais frequência ao estado inicial, o que limita a exploração dos demais estados e torna a convergência para o equilíbrio mais lenta.

    \end{resposta}
\end{enumerate}


\section*{Questão 4: Voltando à origem}
Considere uma cadeia de Markov cujo espaço de estados é um láttice de duas dimensões sobre os números naturais, ou seja, $S = \{(i, j) \mid i \geq 1, j \geq 1\}$. Cada estado pode transicionar para um de seus vizinhos no láttice. Entretanto, se afastar da origem (se movimentar para o norte ou para o leste) tem probabilidade $p/2$, e se aproximar da origem tem probabilidade $(1-p)/2$, onde $p$ é um parâmetro do modelo (nas bordas, utilize self-loops). Assuma que $p \in \{0.25, 0.35, 0.45\}$.

\begin{enumerate}
    \item Construa um simulador para essa cadeia de Markov.
    \begin{resposta}

    \end{resposta}
    \item Utilize o simulador para estimar a distribuição estacionária da origem (estado $(1,1)$), ou seja $\pi_{1,1}$, para cada valor de $p$. Dica: utilize os tempos de retorno!
    \begin{resposta}

    \end{resposta}
    \item Seja $d(t)$ o valor esperado da distância (de Manhattan) entre $X_t$ (o estado no tempo $t$) e a origem. Utilize o simulador para estimar $d(t)$ para $t \in \{10, 100, 1000\}$, para cada valor de $p$. O que você pode concluir?
    \begin{resposta}

    \end{resposta}
\end{enumerate}



\section*{Códigos}

Os códigos utilizados para a resolução dos exercícios estão disponíveis no repositório do GitHub: \url{https://github.com/lhscaldas/CPS767_MCMC/}

% \bibliographystyle{abntex2-num} % Escolha o estilo de citação desejado
% \nocite{sutton2018reinforcement}
% \bibliography{bibliografia} % Nome do arquivo .bib (sem a extensão)

\end{document}